\section*{ЗАКЛЮЧЕНИЕ}
\addcontentsline{toc}{section}{ЗАКЛЮЧЕНИЕ}

В результате выполнения выпускной квалификационной работы была разработана программно-информационная система для организации и проведения занятий йогой — мобильное приложение «Йога-студия» для платформы Android.

В ходе выполнения работы были решены следующие задачи:
\begin{enumerate}
	\item проведён анализ предметной области и существующих аналогов (TimePad, FitStars, ручная запись в Excel);
	\item сформулированы функциональные (FR-01–FR-14) и нефункциональные требования к системе;
	\item обоснован выбор технологий (Kotlin, Android SDK, SharedPreferences, MPAndroidChart, CircleImageView);
	\item разработана архитектура приложения по принципу Single Activity с использованием фрагментов;
	\item спроектирован пользовательский интерфейс, включающий основные экраны и навигацию;
	\item реализованы основные модули: авторизация, каталог занятий, расписание, профиль, запись на занятия, отзывы, статистика;
	\item выполнено тестирование приложения.
\end{enumerate}

Разработанное приложение позволяет пользователям просматривать расписание занятий с фильтрацией по городу, записываться на занятия и отменять запись, просматривать информацию об инструкторах, редактировать профиль, отслеживать прогресс в виде графиков, оставлять отзывы и оценки, а также получать уведомления и напоминания.

Приложение соответствует всем сформулированным требованиям, работает автономно (без подключения к интернету), имеет интуитивно понятный интерфейс на русском языке и поддерживает минимальную версию Android 7.0 (API 24).

Дальнейшее развитие системы может включать добавление серверной части для синхронизации данных между устройствами, интеграцию с платёжными системами для оплаты занятий, расширение аналитики и отчётности для администраторов, а также разработку версии для iOS.